\label{chap:83-17}

\section{Morphismes archimédien et exceptionnel à partir d'un domaine enrichi commun}
\label{p12-83-c17-s1:chap:doble}
\index{double projection}
\index{Moonshine}

La double projection compare deux évaluations du domaine enrichi commun
\(X_\infty^{\rm cal}\) déjà construit dans le
\cref{p12-83-c17-s2:thm:lector-excepcional-global-rev6}. L'une produit des valeurs
archimédiennes ; l'autre, avec un étalonnage de son codomaine, conserve
l'incidence combinatoire qui conduit aux codes, aux plans et aux groupes
d'automorphismes. Les deux évaluations partent de l'unique état initial
\(X_\infty^{\rm cal}\) et préservent la provenance commune dans des codomaines
différenciés.

\subsection{Domaines des \(2\) projections}

Sur \(X_\infty^{\rm cal}\), les cylindres se contractent en un point et la
frontière admet la prolongation d'incidence. Dans cet espace de résidence, nous fixons
\(\mathbf{Arch}:=\mathbb R\). L'évaluation archimédienne et l'évaluation
exceptionnelle sont respectivement
\[
\mathfrak R:X_\infty^{\rm cal}\longrightarrow\mathbf{Arch},
\qquad
\mathfrak E:X_\infty^{\rm cal}\times\mathscr C_{\rm exc}
\longrightarrow\mathbf{Exc}.
\]
Ici, \(\mathscr C_{\rm exc}\) contient l'ordre projectif, la jauge de
Paley, la dodécade, l'alignement d'incidence et la chambre de réseau
spécifiés dans chaque construction. La première application produit des valeurs ; la
seconde, des codes, des plans, des réseaux et des automorphismes. Le théorème global
précédent est la source de la compatibilité et permet de comparer
des invariants obtenus par des opérations distinctes sans intersecter de nouveau
des domaines définis séparément. La comparaison finie de cet ouvrage
s'effectue dans le drapeau
\(B_K\subset H^-_\alpha\).

\begin{figure}[H]
\centering
\resizebox{\textwidth}{!}{\input{figures/fig_doble_proyeccion}}
\caption{Schéma typé des deux projections. L'évaluation archimédienne
identifie des états d'une même fibre ; l'évaluation exceptionnelle, relative
à son étalonnage, conserve la structure d'incidence et les automorphismes.}
\label{p12-83-c17-s1:fig:doble-proyeccion}
\end{figure}

\subsection{Topologie de l'espace des histoires}

La formulation au moyen de cylindres admet une réalisation standard comme
limite projective. Soient \(\mathsf S_n\) les ensembles finis d'états
admissibles à la profondeur \(n\), munis de la topologie discrète, et soient
\[
 \tau_{n+1,n}:\mathsf S_{n+1}\longrightarrow\mathsf S_n
\]
les applications de troncature. Dans la réalisation HMT, ces applications
sont surjectives : la loi de prolongation fournit une continuation
survivante de chaque préfixe, et sa compatibilité projective est conservée à
la limite par le
\cref{p12-83-c17-s2:thm:lector-excepcional-global-rev6}. C'est l'instance
du système qui intervient dans la double projection globale
\eqref{eq:frontal-doble-proyeccion}. Nous définissons
\[
 \mathcal X=\varprojlim_n\mathsf S_n.
\]
Si \(x\ne y\), soit \(m(x,y)\) la première profondeur à laquelle leurs
composantes diffèrent. La fonction
\[
 d(x,y)=3^{-m(x,y)},\qquad d(x,x)=0,
\]
est une ultramétrique qui induit la topologie de la limite projective.

\begin{theorem}[Espace compact des histoires]
\label{p12-83-c17-s1:thm:compacto-historias}
\label{thm:compacto-historias}
L'espace \(\mathcal X\) est compact et totalement discontinu. Lorsque
toute histoire admet des bifurcations à des profondeurs arbitrairement grandes,
\(\mathcal X\) n'a pas de points isolés et est homéomorphe à l'espace de
Cantor.
\end{theorem}

\begin{proof}
Le produit \(\prod_n\mathsf S_n\) est compact. Les équations
\(\tau_{n+1,n}(x_{n+1})=x_n\) définissent un sous-ensemble fermé, de sorte que
\(\mathcal X\) est compact. Les ensembles d'histoires ayant un préfixe fixe
sont à la fois ouverts et fermés et forment une base ; cela prouve la
discontinuité totale. Sous l'hypothèse de bifurcation arbitrairement profonde,
aucun cylindre ne contient d'histoire isolée. La caractérisation des
espaces compacts, métriques, parfaits et totalement discontinus conclut
la dernière affirmation.
\end{proof}

Les intervalles fermés \(I_n(s)\) associés par le lecteur archimédien à
chaque \(s\in\mathsf S_n\) sont compatibles avec la troncature et satisfont
\[
 \operatorname{diam}I_n(s)\leq C\lambda^n,
 \qquad 0<\lambda<1.
\]
L'intersection des intervalles correspondant à une histoire contient
alors un unique point.

\begin{theorem}[Évaluation archimédienne]
\label{p12-83-c17-s1:thm:evaluacion-holder}
\label{thm:evaluacion-holder}
L'application
\[
 \operatorname{val}:\mathcal X\longrightarrow\mathbb R,
 \qquad
 \{\operatorname{val}(x)\}=\bigcap_n I_n(x_n),
\]
est continue et höldérienne par rapport à \(d\). Si
\(x\sim_{\mathrm{Arch}}y\) équivaut à
\(\operatorname{val}(x)=\operatorname{val}(y)\), l'application induite
est un homéomorphisme
\[
 \mathcal X/{\sim_{\mathrm{Arch}}}
 \;\cong\;\operatorname{val}(\mathcal X).
\]
\end{theorem}

\begin{proof}
Deux histoires ayant un préfixe commun de longueur \(n\) ont des valeurs dans un
même intervalle de diamètre inférieur ou égal à \(C\lambda^n\). Puisque
\(d(x,y)=3^{-n}\), cette estimation est une borne de Hölder. L'application
induite sur le quotient est une bijection continue d'un compact sur
un espace séparé et, par conséquent, un homéomorphisme.
\end{proof}

Le théorème formalise la projection archimédienne : la valeur retient la classe
d'évaluation et élimine la distinction entre les histoires situées dans une même
fibre. La couverture cohérente, l'exhaustion par des compacts et la
reconstruction exacte du quotient sont démontrées dans les
\cref{p12-83-c18-s1:thm:sobreyectividad-evaluacion-compacto,p12-83-c18-s1:prop:agotamiento-recta-real-u024,p12-83-c18-s2:hmtch:thm-cociente-real} ;
par conséquent, dans la réalisation HMT active,
\(\operatorname{val}(\mathcal X)=\mathbb R\).

\subsection{Séparation par les \(2\) projections}

\begin{theorem}[Séparation conjointe]
\label{p12-83-c17-s1:thm:separacion-conjunta}
\label{thm:separacion-conjunta}
Soient \(P_1:\mathcal X\to Y_1\) et \(P_2:\mathcal X\to Y_2\) des applications
continues à codomaines séparés. Si
\[
 x\ne y\quad\Longrightarrow\quad
 P_1(x)\ne P_1(y)\ \text{ou}\ P_2(x)\ne P_2(y),
\]
l'application diagonale
\[
 (P_1,P_2):\mathcal X\longrightarrow Y_1\times Y_2
\]
est un plongement topologique.
\end{theorem}

\begin{proof}
L'hypothèse exprime exactement l'injectivité de l'application diagonale.
Une application continue et injective d'un compact dans un espace séparé
est un homéomorphisme sur son image.
\end{proof}

Une jauge déclarée \(c\in\mathscr C_{\rm exc}\) étant fixée, soit
\[
 \mathfrak E_c:X_\infty^{\rm cal}\longrightarrow\mathbf{Exc},
 \qquad
 \mathfrak E_c(x):=\mathfrak E(x,c).
\]
Appliqué à HMT, le théorème prend
\[
 P_1=\mathfrak R,\qquad P_2=\mathfrak E_c
\]
sur le domaine commun, ou bien sur le quotient par les changements de
coordonnées déclarés. La compatibilité projective et la double réalisation
globale sont démontrées respectivement dans les
\cref{p12-83-c17-s2:thm:lector-excepcional-global-rev6,p12-83-c17-s2:thm:doble-realizacion-global-rev6}.
La fidélité est typée par le codomaine : le lecteur exceptionnel est fidèle sur
le flux visible, tandis que le domaine enrichi conserve la mémoire qu'aucune
publication isolée ne contracte. Le passage à toute profondeur est ainsi une
conclusion de la construction, non une hypothèse externe.

À un niveau fini \(S\), cette condition est décidable. Étant donnés
\(p:S\to A\), \(q:S\to B\) et une relation d'équivalence des coordonnées,
on énumère les couples \(x\not\sim y\) qui satisfont simultanément
\[
 p(x)=p(y),\qquad q(x)=q(y).
\]
L'absence de tels couples équivaut à l'injectivité conjointe à ce
niveau. Le calcul doit effectuer les deux applications séparément : lorsque
\(q=f\circ p\), la seconde possède les mêmes fibres que la première et ne
récupère pas l'information éliminée par \(p\).

Au niveau développé ici, l'égalité
\[
 B_K=H_e\cap P_\pi
\]
constitue le témoin fini de compatibilité entre les deux branches. Sa
prolongation compatible à toutes les profondeurs est fournie par les
\cref{p12-83-c17-s2:thm:lector-excepcional-global-rev6,p12-83-c17-s2:thm:doble-realizacion-global-rev6} ;
l'égalité des supports est l'ombre finie de cette réalisation conjointe.

\subsection{Quotient archimédien par les fibres d'évaluation}

Une section HMT contient bloc, orientation, retenue, étiquette d'observable
et coordonnée d'extension.
L'évaluation \(\val\) contracte ses cylindres en un point. Si deux sections
\(x,y\) satisfont
\[
\val(x)=\val(y),
\]
l'évaluation archimédienne les identifie même si leur transport est distinct. En
ce sens précis,
\[
\boxed{\text{une valeur réelle est une classe d’équivalence de l’évaluation}.}
\]
Dans le sous-domaine canonique, somme, produit et évaluation descendent de manière
compatible, et l'image hérite de ces opérations. L'existence de la limite, la
détermination de son image, la caractérisation des fibres et la couverture de
tout \(\RR\) sont établies par les
\cref{p12-83-c18-s1:thm:sobreyectividad-evaluacion-compacto,p12-83-c18-s1:prop:agotamiento-recta-real-u024,p12-83-c18-s2:hmtch:thm-cociente-real}.
Les fenêtres finies présentent la structure locale de ce résultat global ; la
survie unitaire et la compatibilité sont des propriétés démontrées de la
prolongation HMT.

\subsection{Invariant commun avec le système de Witt}

Le vecteur \(K\) définit le support supérieur
\[
B_K=\{m:K_m\ge729\}=\{4,6,9,10\}.
\]
Dans l'atlas de Witt apparaissent les supports
\[
P_\pi=\{2,4,5,6,7,8,9,10,11\},
\]
\[
H_e=\{1,3,4,6,9,10\},
\qquad
H_\varphi=\{1,2,3,4,7,9\}.
\]
Le support des retenues négatives de \(\alpha\) est
\[
H_\alpha^-=\{4,5,6,7,9,10\}.
\]
On vérifie les identités de supports
\[
\boxed{
B_K=H_e\cap P_\pi=H_e\cap H_\alpha^-,}
\]
et
\[
\boxed{H_\alpha^-=B_K\sqcup\{5,7\}.}
\]
La polarité de Witt envoie les quatre points de \(B_K\) sur les paires
\[
\{1,3\},\quad\{2,11\},\quad\{8,12\},\quad\{5,7\}.
\]
Trois reçoivent l'étiquette positive ; la dernière complète le support négatif.
Ces identités sont exactes dans la jauge déclarée et constituent des contrôles
internes d'un même état engendré : le vecteur, le seuil, les mots et la
jauge de Paley sont des coordonnées corrélées de la construction, non des
entrées statistiquement indépendantes. L'égalité s'obtient par
évaluation directe et ne constitue pas une estimation probabiliste.

\subsection{Plan résiduel sur une dodécade et relèvement des hexades}
\label{p12-83-c17-s1:sec:w12-w24-elevacion}

La relation entre les deux plans de Witt s'obtient au sein du code
binaire étendu de Golay et non par une identification de cardinalités.
Soit
\[
\mathcal G_{24}\subset\mathbb F_2^{24}
\]
le code binaire étendu de Golay, et notons \(\mathcal O_{24}\) la famille
des supports de ses mots de poids huit. C'est un résultat classique que
\((\{1,\ldots,24\},\mathcal O_{24})\) est le système de Steiner
\(W_{24}=S(5,8,24)\) \cite{MacWilliamsSloane1977,ConwaySloane1999}.
Fixons un mot de poids douze et écrivons \(D\) pour son support. Nous définissons
\begin{equation}
 \mathcal H(D)
 :=\{O\cap D:O\in\mathcal O_{24},\ |O\cap D|=6\}.
 \label{p12-83-c17-s1:eq:hexadas-residuales-dodecada}
\end{equation}
La famille \(\mathcal H(D)\) est intrinsèque au couple
\((\mathcal G_{24},D)\) : elle n'exige pas de choisir des coordonnées dans \(D\).

\begin{theorem}[Plan résiduel d'une dodécade]
\label{p12-83-c17-s1:thm:dodecada-w12}
Le système d'incidence
\[
 (D,\mathcal H(D))
\]
est un système de Steiner \(S(5,6,12)\). En particulier,
\(|\mathcal H(D)|=132\), et tout sous-ensemble de cinq points de \(D\) est
contenu dans une unique hexade de \(\mathcal H(D)\).
\end{theorem}

\begin{proof}
Soit \(A\subset D\), avec \(|A|=5\). Puisque les octades forment
\(S(5,8,24)\), il existe une unique octade \(O_A\) qui contient \(A\). Si
\(j=|O_A\cap D|\), la somme binaire des mots caractéristiques de
\(O_A\) et \(D\) appartient à \(\mathcal G_{24}\) et a pour poids
\[
 \operatorname{wt}(\mathbf1_{O_A}+\mathbf1_D)=8+12-2j=20-2j.
\]
Les poids du code binaire étendu de Golay sont
\(0,8,12,16,24\). Puisque \(j\leq8\), il s'ensuit que
\(j\in\{2,4,6\}\). L'inclusion \(A\subset O_A\cap D\) donne \(j\geq5\),
donc \(j=6\). Par conséquent, \(O_A\cap D\in\mathcal H(D)\) est une hexade qui
contient \(A\).

Si deux membres de \(\mathcal H(D)\) contenaient \(A\), leurs octades
d'origine contiendraient le même ensemble de cinq points ; la propriété de
Steiner de \(W_{24}\) les ferait coïncider. Cela prouve l'unicité. Le
nombre de blocs s'obtient en comptant les incidences avec des sous-ensembles de cinq
points :
\[
 |\mathcal H(D)|
 =\frac{\binom{12}{5}}{\binom65}=132.\qedhere
\]
\end{proof}

\begin{theorem}[Relèvement unique d'une hexade]
\label{p12-83-c17-s1:thm:elevacion-unica-hexada}
Pour chaque \(H\in\mathcal H(D)\), il existe une unique octade
\(O_H\in\mathcal O_{24}\) telle que
\[
 O_H\cap D=H.
\]
En conséquence, l'application
\[
 \iota_D:\mathcal H(D)\longrightarrow\mathcal O_{24},
 \qquad H\longmapsto O_H,
\]
est bien définie et injective.
\end{theorem}

\begin{proof}
L'existence fait partie de la définition de \(\mathcal H(D)\). Si deux
octades \(O_1,O_2\) avaient pour intersection \(H\) avec \(D\), toutes deux
contiendraient n'importe quel sous-ensemble de cinq points de \(H\). L'unicité de
l'octade incidente à ces cinq points implique \(O_1=O_2\).
L'injectivité découle d'une nouvelle intersection avec \(D\).
\end{proof}

Le relèvement possède une loi locale qui sera utilisée pour comparer les cinq
régions de \(\pi\). Dans un système \(S(t,k,v)\), le nombre de blocs qui
contiennent un sous-ensemble fixé de \(s\leq t\) points est
\[
 \lambda_s=\frac{\binom{v-s}{t-s}}{\binom{k-s}{t-s}}.
\]
Par conséquent,
\begin{equation}
 \lambda_4(W_{12})
 =\frac{\binom81}{\binom21}=4,
 \qquad
 \lambda_4(W_{24})
 =\frac{\binom{20}1}{\binom41}=5.
 \label{p12-83-c17-s1:eq:lambda4-w12-w24}
\end{equation}

\begin{corollary}[Décomposition \(4+1\) sur une face]
\label{p12-83-c17-s1:cor:cuatro-mas-uno-w24}
Soit \(B\subset D\), avec \(|B|=4\). Les cinq octades de \(W_{24}\) qui
contiennent \(B\) se décomposent de manière unique en
\[
 \{O_H:H\in\mathcal H(D),\ B\subset H\}
 \;\sqcup\;\{O_B^\perp\}.
\]
Le premier ensemble contient quatre octades et satisfait
\(|O_H\cap D|=6\) ; l'octade restante satisfait
\[
 O_B^\perp\cap D=B.
\]
\end{corollary}

\begin{proof}
La première égalité de \eqref{p12-83-c17-s1:eq:lambda4-w12-w24} fournit quatre hexades
résiduelles contenant \(B\), et le
\cref{p12-83-c17-s1:thm:elevacion-unica-hexada} fournit quatre octades distinctes. La
seconde égalité de \eqref{p12-83-c17-s1:eq:lambda4-w12-w24} montre qu'il existe exactement
une cinquième octade \(O_B^\perp\).

Puisque \(B\subset O_B^\perp\cap D\), le spectre des intersections employé dans
la preuve du \cref{p12-83-c17-s1:thm:dodecada-w12} ne laisse que les possibilités quatre et
six. Une intersection de taille six serait une cinquième hexade résiduelle qui
contient \(B\), ce qui contredit \(\lambda_4(W_{12})=4\). L'intersection
est donc de taille quatre et coïncide avec \(B\).
\end{proof}

\subsubsection{Alignement avec les \(12\) positions HMT}

Le plan résiduel est défini sur les points de \(D\), tandis que le
plan obtenu à partir du code ternaire \(C_W\) est défini sur les douze
positions dodécaphasiques. Pour ne pas confondre les deux ensembles, nous appellerons
\emph{alignement binaire HMT} une bijection
\begin{equation}
 \ell:\{1,\ldots,12\}_{\rm HMT}\xrightarrow{\sim}D
 \quad\text{tel que}\quad
 \{\ell(H):H\in W_{12}^{\rm HMT}\}=\mathcal H(D).
 \label{p12-83-c17-s1:eq:alineamiento-hmt-w24}
\end{equation}
Ici, \(W_{12}^{\rm HMT}\) désigne le plan des supports de poids six
obtenu à partir du code ternaire \(C_W\). L'unicité du système de Witt
\(S(5,6,12)\) à isomorphisme près garantit l'existence de \(\ell\). Deux
alignements distincts diffèrent par des automorphismes des plans de départ et
d'arrivée ; c'est pourquoi \(\ell\) est une donnée de coordonnées et non un invariant
supplémentaire.
Le couple \((D,\ell)\) est la donnée d'étalonnage nécessaire pour transporter
le relèvement \(\iota_D\) vers le système de coordonnées HMT\@. Une fois celui-ci fixé, chaque
hexade HMT \(H\) détermine sans ambiguïté l'octade
\[
 \widehat H:=\iota_D(\ell(H)).
\]
La dodécade et l'alignement sont des données du codomaine binaire ; ils ne s'identifient pas
à une sortie décimale du noyau de transport.

Pour la face distinguée
\[
 B_K=\{4,6,9,10\},
\]
les quatre hexades incidentes sont
\[
 B_K\sqcup\{1,3\},\quad
 B_K\sqcup\{2,11\},\quad
 B_K\sqcup\{5,7\},\quad
 B_K\sqcup\{8,12\}.
\]
La troisième coïncide avec
\(H_\alpha^-=B_K\sqcup\{5,7\}\). L'alignement \(\ell\) relève ces
quatre hexades en les quatre octades résiduelles sur \(\ell(B_K)\) ; le
\cref{p12-83-c17-s1:cor:cuatro-mas-uno-w24} détermine en outre une unique octade transversale.

\subsubsection{Correspondance d'incidence entre la microfibre de \texorpdfstring{\(\pi\)}{pi} et les octades résiduelles : multiplicité \texorpdfstring{\(432+4\cdot144=1008\)}{432+4x144=1008}}

La microfibre exacte du mot
\(w_\pi=010211\) contient cinq régions décimales distinctes :
\[
\begin{array}{c|c|c}
\text{multiplicité}&U_6& e(U_6)\\
\hline
432&501|614|498|169|272|272&(5,5,5,5,5,5)\\
144&810|923|870|169|272|272&(3,3,9,5,5,5)\\
144&810|983|810|169|272|272&(3,9,3,5,5,5)\\
144&870|923|810|169|272|272&(9,3,3,5,5,5)\\
144&870|923|810|418|674|521&(9,3,3,7,1,7).
\end{array}
\]
Ici, \(e(U_6)\) est la signature multiplicative récupérée à partir de chaque bloc
décimal \(\overline{d_1d_2d_3}\) au moyen de
\(e=d_2-d_1\pmod {10}\). La première région est l'unique région de multiplicité
432 et de signature constante. Les quatre autres sont de multiplicité 144 ; parmi
elles, trois partagent le bloc de retour \(169|272|272\), tandis que la
quatrième possède le retour \(418|674|521\). Le catalogue distingue les cinq
lignes au moyen d'un prédicat reproductible : les trois lignes de type
\(\mathrm{ES}\) forment le triplet positif ; parmi les deux lignes de type
\(\mathrm{NO}\), celle de multiplicité maximale et de signature constante est appelée
transversale, et l'autre est appelée négative. Relativement à ce prédicat de
rôle, la partition est
\[
 \boxed{1_\perp+3_{++}+1_{--}.}
\]
Les noms de signe sont une convention ; ce qui est intrinsèque au catalogue est la
partition obtenue à partir de ses colonnes, non une polarité binaire préalable.

L'alignement régional complète la donnée \((D,\ell)\) par une bijection
qui envoie la région transversale sur \(O_{\ell(B_K)}^\perp\), la région
négative sur \(\widehat{H_\alpha^-}\) et les trois régions positives sur les
relèvements de
\[
 B_K\sqcup\{1,3\},\qquad
 B_K\sqcup\{2,11\},\qquad
 B_K\sqcup\{8,12\}.
\]
L'affectation ordonnée entre les trois régions positives et ces trois hexades
est un choix de coordonnées. Les deux lignes distinguées et le triplet sont
reproductibles au moyen du prédicat de catalogue précédent ; leur dénomination comme
négative, transversale et positives relève de l'alignement. Ainsi, la loi \(4+1\) n'identifie pas par la seule
cardinalité les cinq régions à cinq octades : elle fournit l'espace
d'incidence, tandis que \((D,\ell)\) et l'alignement régional spécifient
le transport entre les deux objets.

\subsection{Prolongation en réseau et écrans dimensionnels}

Le relèvement des hexades ouvre deux branches qu'il ne faut pas confondre. La première
conserve l'incidence binaire et prolonge \(W_{12}\) vers \(W_{24}\). La
seconde conserve le transport ternaire :
\[
C_{12}^{(3)}
\longrightarrow N(A_2^{12})
\longrightarrow\Lambda_{24}
\longrightarrow V^\natural
\longrightarrow\mathbb M
\longrightarrow\text{Moonshine}.
\]
La preuve primaire du recollement, du \(3\)-voisin sans racines et du corridor
jusqu'à Moonshine réside dans le
\cref{chap:bandera-excepcional-acumulativa,thm:corredor-hmt-moonshine}.
Il n'existe ni flèche canonique \(W_{24}\to N(A_2^{12})\), ni action
directe du Monstre sur les mots TPK : les deux branches partent du même
code étalonné et préservent des structures distinctes.

Les réductions \(243/11\), \(12\to11\to10\) et \(26\to24\) ne sont pas non plus des
identités entre numéraux. Le \cref{chap:pantallas-dualidad} construit leurs
espaces, quotients et applications : le quotient parfait ternaire, l'écran
marqué par \(K\) et la réduction isotrope du réseau lorentzien. Cette
section conserve ainsi une seule fonction : relier l'incidence finie déjà
démontrée à ses deux prolongations, sans répéter les preuves de référence.

\subsection{Involution résidu–quotient et conservation de l'information de l'état}

Pour \(n=9Q+r\), nous définissons l'énergie de réseau
\[
E_R(r,Q)=\frac{r^2}{R^2}+Q^2R^2.
\]
Alors
\[
\boxed{E_R(r,Q)=E_{1/R}(Q,r).}
\]
L'échange \((r,Q,R)\leftrightarrow(Q,r,1/R)\) est une involution exacte
du modèle de réseau : le résidu et le quotient échangent leurs fonctions.

\subsection{Compatibilité finie et prolongation conjointe}

Le témoin fini \(B_K\), le motif de retenues et le cadre de Witt fixent la
compatibilité locale. Sa prolongation n'est plus laissée à l'état de schéma
abstrait : le \cref{p12-83-c17-s2:thm:lector-excepcional-global-rev6} construit bloc par
bloc le lecteur exceptionnel avec une unique jauge initiale transportée, et le
\cref{p12-83-c17-s2:thm:doble-realizacion-global-rev6} démontre que cette lecture et
l'évaluation archimédienne se globalisent sur la même section TPK. Tronquer
l'histoire avant de lire équivaut dans les deux cas à tronquer la sortie après
lecture.

L'évaluation archimédienne identifie les histoires par leur valeur limite.
L'évaluation exceptionnelle conserve le flux visible et son incidence
Paley--Witt ; elle ne récupère pas pour autant les mémoires cachées que l'émetteur visible
ne publie pas. La double projection est donc une construction conjointe avec
deux codomaines distincts, non une identification entre le continu et la branche
exceptionnelle.

\begin{remark}[Compatibilité des deux projections]
Le support \(B_K\), le motif de retenues et le cadre de Witt fournissent les
invariants communs. La chaîne de réseaux prolonge cette compatibilité jusqu'à un
voisin sans racines et, par classification classique, jusqu'au réseau de Leech. La
branche binaire vers \(W_{24}\) demeure séparée : son morphisme d'incidence est
\(\iota_D\), défini après fixation de la dodécade et de l'alignement
\((D,\ell)\), et n'a pour codomaine ni un réseau de Niemeier ni le réseau
de Leech.
\end{remark}


\section{Projection exceptionnelle de l'état enrichi : incidence compatible à toute profondeur}
\label{p12-83-c17-s2:chap:lector-excepcional-global-rev6}
\index{lecteur exceptionnel}\index{double projection}
\index{Paley--Witt}\index{système projectif}

Le témoin fini de la double projection acquiert une portée globale lorsque la
lecture exceptionnelle est appliquée à chaque bloc visible et que la jauge est
transportée, au lieu d'être choisie de nouveau, à mesure que la profondeur augmente. Cette
section construit cette famille. Le domaine demeure l'histoire TPK
enrichie ; la lecture exceptionnelle est fidèle sur son flux visible et ne
prétend pas se substituer aux champs de mémoire que ce flux ne publie pas.

\subsection{L'opération locale sur les 729 mots}

Dans la carte de Paley--Witt, on fixe
\[
A_W=
\begin{pmatrix}
0&1&1&1&1&1\\
1&0&1&2&2&1\\
1&1&0&1&2&2\\
1&2&1&0&1&2\\
1&2&2&1&0&1\\
1&1&2&2&1&0
\end{pmatrix}
\in M_6(\mathbb F_3).
\]
La multiplication matricielle donne
\[
A_W^{\mathsf T}=A_W,\qquad A_W^2=-I_6.
\]
Pour tout opérateur conjugué \(A=fA_Wf^{-1}\), nous définissons
\[
\Gamma_A:\mathbb F_3^6\longrightarrow\mathbb F_3^{12},
\qquad
\Gamma_A(w)=(w,wA).
\]
Son image est une copie étalonnée \(C_W(A)\) du code ternaire de Witt.
La projection sur les six premières coordonnées est un inverse à gauche :
\[
\operatorname{pr}_1\Gamma_A=\operatorname{id}_{\mathbb F_3^6}.
\]
Par conséquent, \(\Gamma_A\) est totale et injective sur les \(729\) mots et
n'a pas besoin de distinguer à l'avance \(\pi\), \(e\), \(\varphi\) ni aucune
autre sortie.

\subsection{Prolongations compatibles et unicité de la jauge initiale}

Soit \(X_n\) l'ensemble des préfixes TPK admissibles de profondeur \(n\), avec
troncatures
\[
p_{n,m}:X_n\longrightarrow X_m,\qquad m\leq n.
\]
Fixons un cadre exceptionnel initial \(f_0\). Puisque les changements \(g_j\) font
partie du transport enregistré par l'histoire, un préfixe étalonné est le
couple \((x,f_0)\) ; nous écrivons
\[
 X_n^{\rm cal}:=X_n\times\{f_0\},\qquad
 p_{n,m}^{\rm cal}(x,f_0):=(p_{n,m}x,f_0),
\]
et
\[
 X_\infty^{\rm cal}
 :=\varprojlim_n(X_n^{\rm cal},p_{n,m}^{\rm cal}).
\]
Dans la suite, nous abrégeons \(p_{n,m}^{\rm cal}\) en \(p_{n,m}\).
L'émission visible compatible est
\[
\varepsilon_n:X_n^{\rm cal}\longrightarrow(\mathbb F_3^6)^n,
\qquad
\varepsilon_n(x)=(w_0,\ldots,w_{n-1}),
\]
et satisfait
\[
\operatorname{pr}_{n,m}\varepsilon_n
=\varepsilon_m p_{n,m}.
\]
Cette application distingue ses trois types :
\[
|\mathcal U_6|=468,\qquad
|\operatorname{im}w_6|=243,\qquad
|\mathbb F_3^6|=729.
\]
L'émission \(w_j\) ne s'identifie pas à l'état complet qui conserve
feuillet, quotient, retenue, orientation, phase et survie.

La jauge exceptionnelle se compose d'une carte de Witt, d'une dodécade, d'un
alignement, d'une origine de réseau et d'une chambre orientée. Elle est fixée une seule
fois dans le bloc initial. Pour
\(x\in X_n^{\rm cal}\), le transport de l'histoire fournit
\[
 g_j(x)\in\operatorname{GL}_6(\mathbb F_3),
 \qquad 0\le j<n-1,
\]
et \(g_j(x)\) ne dépend que du préfixe
\(p_{n,j+1}(x)\). Le cadre est transporté par
\[
 f_{j+1}(x)=g_j(x)f_j(x),\qquad
 A_j=f_jA_Wf_j^{-1}.
\]
Ainsi, la famille \((A_j)\) est déterminée par l'histoire et par la jauge
initiale ; ce n'est pas une suite de choix indépendants.

Soit \(\mathcal F_{\rm cal}\subseteq\operatorname{GL}_6(\mathbb F_3)\)
l'ensemble des cadres accessibles à partir de \(f_0\). L'émission qui conserve la
jauge transportée est
\[
 \widetilde\varepsilon_n:X_n^{\rm cal}
 \longrightarrow(\mathbb F_3^6\times\mathcal F_{\rm cal})^n,
 \qquad
 \widetilde\varepsilon_n(x)
 =\bigl((w_0,f_0),\ldots,(w_{n-1},f_{n-1})\bigr).
\]
La dépendance par préfixes implique
\(\operatorname{pr}_{n,m}\widetilde\varepsilon_n
=\widetilde\varepsilon_m p_{n,m}\). Sa projection sur les mots est
\(\varepsilon_n\) ; l'oubli des cadres ne suffit pas à reconstruire la
lecture exceptionnelle.

\subsection{Naturalité et compatibilité à toute profondeur}

L'opération locale est covariante avec la convention des vecteurs-lignes. Pour
tout changement de coordonnées \(f\in\operatorname{GL}_6(\mathbb F_3)\),
\[
A^f=fAf^{-1}
\quad\Longrightarrow\quad
\Gamma_{A^f}(wf^{-1})
=\Gamma_A(w)(f^{-1}\oplus f^{-1}).
\]
Le changement de cadre transporte simultanément le mot, l'opérateur et
l'atlas. Il n'est pas exigé qu'il demeure dans le stabilisateur d'une unique copie
étiquetée.

L'ensemble des cadres accessibles \(\mathcal F_{\rm cal}\) étant déjà défini,
nous construisons l'atlas étalonné fixe
\[
 \mathscr C_W^{\rm cal}
 :=
 \left\{
 \bigl(f,\Gamma_{fA_Wf^{-1}}(w)\bigr):
 f\in\mathcal F_{\rm cal},\ w\in\mathbb F_3^6
 \right\}.
\]
L'étiquette \(f\) conserve la copie étalonnée dans laquelle vit chaque mot.
Nous définissons
\[
Q_n(x)=
\bigl(
\bigl(f_0,\Gamma_{A_0}(w_0)\bigr),\ldots,
\bigl(f_{n-1},\Gamma_{A_{n-1}}(w_{n-1})\bigr)
\bigr)\in\bigl(\mathscr C_W^{\rm cal}\bigr)^n.
\]
Soit
\[
 r_{n,m}:
 \bigl(\mathscr C_W^{\rm cal}\bigr)^n
 \longrightarrow
 \bigl(\mathscr C_W^{\rm cal}\bigr)^m
\]
la troncature des \(n-m\) derniers blocs.

\begin{theorem}[Compatibilité projective de l'incidence exceptionnelle]
\label{p12-83-c17-s2:thm:lector-excepcional-global-rev6}
\label{thm:lector-excepcional-global-rev6}
Pour tout \(m\leq n\),
\[
\boxed{r_{n,m}\circ Q_n=Q_m\circ p_{n,m}.}
\]
En conséquence, il existe une unique application
\[
\boxed{
Q_\infty:X_\infty^{\rm cal}
\longrightarrow
\bigl(\mathscr C_W^{\rm cal}\bigr)^{\mathbb N}}
\]
dont les projections finies sont les \(Q_n\). L'application est fidèle sur le
flux visible.
\end{theorem}

\begin{proof}
Écrivons
\[
\varepsilon_n(x)
=(w_0,\ldots,w_{m-1},w_m,\ldots,w_{n-1}).
\]
La compatibilité de \(\varepsilon\) montre que le flux visible de
\(p_{n,m}(x)\) est \((w_0,\ldots,w_{m-1})\). Pour \(j<m\), le cadre
\(f_j\) ne dépend que de la jauge initiale et des changements antérieurs à
\(j\), tous contenus dans ce préfixe. Les \(m\) premiers opérateurs
\(A_j\) sont donc les mêmes avant et après troncature. Appliquer
\(\Gamma_{A_j}\) bloc par bloc donne l'identité de l'énoncé. La propriété
universelle de la limite projective produit \(Q_\infty\). Enfin,
\(\operatorname{pr}_1\Gamma_{A_j}=I\) récupère chaque \(w_j\), ce qui prouve
la fidélité sur le flux visible.
\end{proof}

La dernière affirmation conserve son type. Deux histoires ayant le même flux
\((w_j)_j\) et des mémoires cachées distinctes peuvent avoir la même image
exceptionnelle. Le lecteur global préserve exactement ce que publie chaque bloc
visible et l'incidence ajoutée à ce bloc.

La compatibilité des émissions enrichies induit une unique application
\[
 \widetilde\varepsilon_\infty:X_\infty^{\rm cal}
 \longrightarrow
 (\mathbb F_3^6\times\mathcal F_{\rm cal})^{\mathbb N}.
\]
Soient \(\operatorname{pr}_w\) la projection sur le flux de mots et
\[
 \widetilde\Gamma_\infty
 \bigl((w_j,f_j)_{j\ge0}\bigr)
 :=\bigl(
   (f_j,\Gamma_{f_jA_Wf_j^{-1}}(w_j))
 \bigr)_{j\ge0}.
\]
Alors
\(Q_\infty=\widetilde\Gamma_\infty
\circ\widetilde\varepsilon_\infty\), avec domaines et codomaines déclarés.

\subsection{Bifurcation d'un préfixe survivant en valeur archimédienne et incidence exceptionnelle}

\(m\in\mathbb Z\) étant fixé, notons
\[
 \operatorname{flat}:(\mathbb F_3^6)^{\mathbb N}
 \longrightarrow\mathbb F_3^{\mathbb N}
\]
la concaténation des blocs dans l'ordre des coordonnées déclaré, et
\[
 \operatorname{ev}_3((d_k)_{k\ge1})
 :=\sum_{k\ge1}\frac{d_k}{3^k}\in[0,1],
 \qquad
 \tau_m(t):=m+t,
\]
l'évaluation ternaire et la translation par la partie entière. Le préfixe visible
\((d_1,\ldots,d_{6n})\) détermine le cylindre
\[
I_n(x)=
\left[
m+\sum_{k=1}^{6n}\frac{d_k}{3^k},
m+\sum_{k=1}^{6n}\frac{d_k}{3^k}+3^{-6n}
\right].
\]
Une prolongation restreint le cylindre et
\(\operatorname{diam}I_n=3^{-6n}\to0\). Par conséquent, toute histoire infinie
compatible détermine un unique point réel
\(\mathcal P_{\rm arq}(x)\).

\begin{theorem}[Double réalisation globale]
\label{p12-83-c17-s2:thm:doble-realizacion-global-rev6}
\label{thm:doble-realizacion-global-rev6}
Toute section compatible \(x=(x_n)_n\in X_\infty^{\rm cal}\) détermine,
à partir des mêmes préfixes, deux sorties compatibles :
\[
\mathcal P_{\rm arq}(x)\in\mathbb R,
\qquad
\mathcal P_{\rm exc}(x)=Q_\infty(x)
\in\bigl(\mathscr C_W^{\rm cal}\bigr)^{\mathbb N}.
\]
Tronquer avant de lire équivaut à tronquer chaque lecture après lecture. La
première sortie identifie les histoires par leur valeur limite ; la seconde conserve
chaque bloc visible et lui adjoint sa rotation de Paley--Witt.
\end{theorem}

\begin{proof}
L'emboîtement des cylindres et la convergence de leurs diamètres prouvent la
première affirmation. Le
\cref{p12-83-c17-s2:thm:lector-excepcional-global-rev6} prouve la seconde et la
compatibilité avec les troncatures. Toutes deux se factorisent par la même émission :
\[
\mathcal P_{\rm arq}
=\tau_m\circ\operatorname{ev}_3\circ\operatorname{flat}
 \circ\operatorname{pr}_w
 \circ\widetilde\varepsilon_\infty,
\qquad
\mathcal P_{\rm exc}
=\widetilde\Gamma_\infty\circ\widetilde\varepsilon_\infty.
\]
\end{proof}

La double projection est ainsi construite sur des histoires admissibles et une
jauge initiale transportée. Le continu est l'évaluation archimédienne de
la section ; la branche exceptionnelle est la réalisation compatible de son flux
visible. À partir de \(C_W\) s'ouvrent, avec les alignements déjà déclarés, les
branches vers \(W_{12}\), \(W_{24}\) et vers
\(N(A_2^{12})\), Leech, \(V^\natural\) et le Monstre. Ces branches ne transforment pas
les chiffres en un ensemble sur lequel le Monstre agirait directement :
elles transportent l'incidence que l'évaluation par valeur cesse de voir.

\subsection{Couplage postgénératif avec la clôture opératorielle d'Euler}
\label{p12-83-c17-s2:sec:acoplamiento-euler-doble-proyeccion-rev3}

Soit \(x_\pi\in X_\infty^{\rm cal}\) la section survivante du canal de
\(\pi\). La double réalisation du
\cref{p12-83-c17-s2:thm:doble-realizacion-global-rev6} produit
\[
 \mathcal P_{\rm arq}(x_\pi)\in\mathbb R,
 \qquad
 \mathcal P_{\rm exc}(x_\pi)=Q_\infty(x_\pi).
\]
La première branche publie l'angle engendré ; la seconde conserve le flux
visible et réalise dans chaque bloc la relation quadratique de Paley--Witt.

\begin{theorem}[Couplage postgénératif Euler--double projection]
\label{p12-83-c17-s2:thm:acoplamiento-euler-doble-proyeccion-rev3}
Supposons que la construction coinductive du canal de \(\pi\) détermine
\[
 \mathcal P_{\rm arq}(x_\pi)=\pi_{\rm HMT},
\]
et que la branche exceptionnelle utilise la jauge transportée avec
\(A_W^2=-I_6\). Soit \(J_{\mathrm e}\) la réalisation réelle de la source
entière du \cref{thm:dos-realizaciones-fuente-cuadratica-rev3}. Alors
\begin{equation}
 \boxed{
 \operatorname{Exp}\!\bigl(
   \mathcal P_{\rm arq}(x_\pi)J_{\mathrm e}
 \bigr)+I=0.}
 \label{p12-83-c17-s2:eq:acoplamiento-euler-doble-proyeccion-rev3}
\end{equation}
\end{theorem}

\begin{proof}
D'après le \cref{thm:dos-realizaciones-fuente-cuadratica-rev3}, \(A_W\) et
\(J_{\mathrm e}\) réalisent la relation \(u^2+1=0\) dans des catégories distinctes.
Dans la branche réelle,
\[
 \operatorname{Exp}(\theta J_{\mathrm e})
 =\cos\theta\,I+\sin\theta\,J_{\mathrm e}.
\]
La substitution de la sortie précédemment engendrée
\(\theta=\mathcal P_{\rm arq}(x_\pi)=\pi_{\rm HMT}\) produit \(-I\), et
l'addition de \(I\) donne l'identité annoncée.
\end{proof}

\paragraph{Ordre causal et absence de circularité.}
La dépendance démontrée est
\[
 \text{émetteur}
 \longrightarrow x_\pi
 \longrightarrow\mathcal P_{\rm arq}(x_\pi)
 \longrightarrow\pi_{\rm HMT}
 \longrightarrow
 \operatorname{Exp}(\pi_{\rm HMT}J_{\mathrm e})+I.
\]
La branche exceptionnelle fixe le type quadratique ; elle ne sélectionne ni l'angle ni ses
chiffres. L'identité d'Euler reconnaît ensuite la compatibilité des deux
réalisations et n'intervient pas comme entrée du générateur de
\(\pi_{\rm HMT}\).


\section{Globalisation naturelle des projections archimédienne et exceptionnelle sur un système projectif commun}
\label{p12-83-c17-s3:chap:globalizacion-doble-proyeccion}

La thèse de la double projection acquiert un contenu mathématique lorsque les deux
lectures sont définies sur un même système d'états et commutent avec la
restriction de profondeur. L'évaluation archimédienne contracte une famille de
cylindres en une valeur ; l'évaluation exceptionnelle conserve des relations
d'incidence que la première peut oublier. Il ne s'agit pas de deux expressions
numériques juxtaposées, mais de deux transformations naturelles de domaine
commun.

\subsection{Systèmes compatibles à profondeur finie}

Soient \((X_n,p_{n+1,n})\), \((A_n,a_{n+1,n})\) et
\((E_n,e_{n+1,n})\) trois systèmes projectifs. Le premier contient des histoires
TPK ; le deuxième, des cylindres ou approximants archimédiens ; le troisième, des données
d'incidence exceptionnelle. La construction active fournit les applications
\[
 R_n:X_n\longrightarrow A_n,
 \qquad
 Q_n:X_n\times C_n\longrightarrow E_n,
\]
où \(C_n\) réunit les données de coordonnées du lecteur exceptionnel. Les
jauges forment une famille compatible et, pour tout \(n\), les diagrammes
commutent
\[
 a_{n+1,n}\circ R_{n+1}=R_n\circ p_{n+1,n},
\]
\[
 e_{n+1,n}\circ Q_{n+1}
 =Q_n\circ(p_{n+1,n}\times c_{n+1,n}).
\]

\begin{theorem}[Globalisation des deux évaluations]
\label{p12-83-c17-s3:thm:globalizacion-evaluaciones}
D'après les compatibilités précédentes, déjà construites bloc par bloc, il existe
des applications canoniques
\[
 R_\infty:\varprojlim X_n\longrightarrow\varprojlim A_n,
 \qquad
 Q_\infty:\varprojlim X_n\times\varprojlim C_n
            \longrightarrow\varprojlim E_n.
\]
Leur application diagonale
\[
 (R_\infty,Q_\infty):X_\infty\times C_\infty
 \longrightarrow A_\infty\times E_\infty
\]
est unique parmi les applications dont les projections coordonnées sont
\(R_n\) et \(Q_n\).
\end{theorem}

\begin{proof}
Soit \(x=(x_n)_n\in\varprojlim X_n\). La première équation de commutativité
implique
\[
 a_{n+1,n}(R_{n+1}(x_{n+1}))=R_n(x_n),
\]
de sorte que \((R_n(x_n))_n\) appartient à \(\varprojlim A_n\). Nous définissons
\(R_\infty(x)\) au moyen de cette famille. Le même argument, appliqué à
\((x_n,c_n)_n\), définit \(Q_\infty\). L'unicité découle du fait qu'un élément
d'une limite projective est déterminé par toutes ses coordonnées.
\end{proof}

Le théorème consigne la preuve concrète réalisée dans HMT : le résultat
archimédien et le code exceptionnel proviennent des mêmes états
compatibles et respectent les applications de troncature.

\subsection{Contraction archimédienne de cylindres compatibles et réalisation de la droite réelle}

À chaque \(x_n\in X_n\) est associé un intervalle compact non vide
\(I_n(x_n)\subset\mathbb R\). Si
\[
 p_{n+1,n}(x_{n+1})=x_n
 \quad\Longrightarrow\quad
 I_{n+1}(x_{n+1})\subseteq I_n(x_n)
\]
et s'il existe \(\varepsilon_n\to0\) avec
\(\operatorname{diam}I_n(x_n)\leq\varepsilon_n\), l'intersection
\[
 \bigcap_{n\geq0}I_n(x_n)
\]
contient un unique réel. On obtient ainsi
\[
 \operatorname{ev}_{\rm arq}:X_\infty^{\rm arq}\longrightarrow\mathbb R.
\]
La phrase « les cylindres tendent vers zéro » signifie ici que leur diamètre
converge vers zéro, non que le nombre évalué tende vers zéro.

La relation
\[
 x\sim_{\rm arq}y
 \quad\Longleftrightarrow\quad
 \operatorname{ev}_{\rm arq}(x)=\operatorname{ev}_{\rm arq}(y)
\]
décrit exactement l'information éliminée par le lecteur. La valeur réelle
est une classe d'évaluation ; l'histoire conserve en outre orientation,
retenue, feuillet et signature de frontière.

\subsection{Projection exceptionnelle par conservation de l'incidence et de la frontière}

La seconde évaluation ne prend pas ses valeurs dans \(\mathbb R\). Son codomaine est une
catégorie de configurations d'incidence. Dans la réalisation finie en vigueur,
la donnée de coordonnées contient :
\[
 \begin{aligned}
 C_{\rm exc}=(\,&\text{ordre projectif},\text{calibre de Paley},
 \text{dodécade},\\
 &\text{alignement},\text{origine réticulaire},\text{chambre}).
 \end{aligned}
\]
Après fixation de cette donnée, la chaîne contient deux descendances qui doivent
demeurer séparées :
\[
 C_W\longrightarrow W_{12}\longrightarrow W_{24},
\]
\[
 C_W\longrightarrow N(A_2^{12})
 \longrightarrow\Lambda_{24}
 \longrightarrow V^\natural
 \longrightarrow\mathbb M.
\]
La première transporte des plans ; la seconde conduit, au moyen de résultats
classiques ultérieurs, au réseau de Leech, à l'algèbre d'opérateurs de
sommets et au groupe Monstre. L'holonomie centrale et le corridor
Schläfli--\(E_6\) forment une extraction parallèle de frontière ; ils ne s'obtiennent pas
en remplaçant l'une des deux descendances précédentes.

\subsection{Compatibilité finie du domaine commun}

Le sceau dodécaphasique
\[
 K=(234,543,140,729,659,824,621,58,914,794,146,601)
\]
détermine le support
\[
 B_K=\{m:K_m\geq729\}=\{4,6,9,10\}.
\]
Dans la carte de Witt publiée,
\[
 P_\pi=\{2,4,5,6,7,8,9,10,11\},
 \qquad
 H_e=\{1,3,4,6,9,10\},
\]
et, par évaluation indépendante des deux supports,
\[
 \boxed{B_K=H_e\cap P_\pi.}
\]
En outre, pour le support de retenue négative
\[
 H_\alpha^-=\{4,5,6,7,9,10\}
\]
on a
\[
 B_K=H_e\cap H_\alpha^-,
 \qquad H_\alpha^-=B_K\sqcup\{5,7\}.
\]
Ces identités constituent un témoin de compatibilité finie : le même
sous-ensemble est retrouvé à partir du sceau arithmétique et de l'incidence de
Witt. Ce n'est pas une égalité de cardinalités, puisqu'elle identifie des positions
concrètes.

\subsection{Séparation typée des publications archimédienne et exceptionnelle de l'état commun}

\begin{theorem}[Critère de fidélité conjointe]
Soient \(X\) compact et \(Y,Z\) des espaces séparés. Si deux applications
continues \(R:X\to Y\) et \(Q:X\to Z\) satisfont
\[
 x\neq y\Longrightarrow R(x)\neq R(y)\ \text{ou}\ Q(x)\neq Q(y),
\]
alors \((R,Q):X\to Y\times Z\) est un plongement topologique.
\end{theorem}

\begin{proof}
La condition est l'injectivité de l'application diagonale. Toute application
continue et injective d'un compact dans un espace séparé est un
homéomorphisme sur son image.
\end{proof}

À chaque profondeur finie, la condition se décide par énumération des
couples \(x\neq y\) qui demeurent simultanément dans une fibre de \(R_n\) et
de \(Q_n\). La compatibilité de ces contrôles lors du passage de \(n+1\) à \(n\)
est précisément l'identité projective démontrée dans le
\cref{p12-83-c17-s2:thm:lector-excepcional-global-rev6}, et le
\cref{p12-83-c17-s2:thm:doble-realizacion-global-rev6} effectue sa
globalisation. Le témoin \(B_K\) présente l'invariant fini commun ;
l'obstruction du lecteur minimal du
\cref{p12-83-c17-s4:thm:obstruccion-lector-excepcional-minimo} délimite le
quotient visible, tandis que le lecteur enrichi conserve les données
d'orientation, de feuillet, de signature, de retour et de frontière requises par la double
publication.

\subsection{Formulation précise de la thèse sur la droite réelle}

L'objet enrichi associé à la lecture archimédienne est la flèche
\[
 \mathbb R_{\mathfrak H}:=
 \bigl(X_\infty^{\rm arq}
 \xrightarrow{\operatorname{ev}_{\rm arq}}
 \operatorname{Im}(\operatorname{ev}_{\rm arq})\subseteq\mathbb R\bigr).
\]
Le domaine conserve les généalogies ; l'image conserve les valeurs. La couverture des
cylindres à toute échelle et la prolongeabilité compatible sont démontrées dans
la construction du lecteur archimédien
\eqref{eq:continuo-lector-arquimediano} : pour tout compact
\(K\subset\mathbb R\),
\(\operatorname{ev}_{\mathbb R}(\Omega_{\infty,K})=K\), et la réunion dirigée
de ces intervalles couvre \(\mathbb R\). Par conséquent,
\[
 \operatorname{Im}(\operatorname{ev}_{\rm arq})=\mathbb R,
 \qquad
 X_\infty^{\rm arq}/\!\sim_{\mathbb R}\ \cong\mathbb R.
\]
Le corps réel n'agit pas comme point de départ de la construction, mais comme
quotient archimédien d'un espace de sections avec mémoire. La descente de la
somme et du produit est vérifiée opération par opération et ne conditionne pas cette
surjectivité déjà démontrée.

La double lecture s'exprime alors par
\[
 X_\infty^{\rm com}\times C_\infty
 \xrightarrow{\ (\operatorname{ev}_{\rm arq},Q_\infty)\ }
 \mathbb R\times E_\infty.
\]
Cette formule est la traduction mathématique de l'affirmation conceptuelle : une
même histoire peut se contracter jusqu'à une valeur et, dans la lecture transversale
de sa frontière, conserver l'incidence qui conduit aux symétries
exceptionnelles.


% Fuente normativa activa de la obstrucción de factorización.
\section{Dépendance du lecteur d'incidence à l'égard de l'état enrichi}
\label{p12-83-c17-s4:sec:obstruccion-lector-excepcional-minimo}

La compatibilité finie des deux évaluations n'implique pas que n'importe quel
lecteur d'incidence raffine la partition produite par le lecteur visible. Le
catalogue régional permet de trancher cette question exactement pour la
réalisation minimale de Paley--Witt. Soit \(\mathcal U_{468}\) l'ensemble
des 468 états régionaux publiés et soit
\[
 \Psi:\mathcal U_{468}\longrightarrow\mathcal W_{243}\subset\mathbb F_3^6,
 \qquad
 \Psi(U_1,\ldots,U_6)=(-U_1,\ldots,-U_6)\pmod 3.
\]
Pour \(w\in\mathbb F_3^6\), l’application
\[
 c(w)=(w,wA_W)\in C_W\subset\mathbb F_3^{12}
\]
est le codeur ternaire fixé dans la section de Paley--Witt. Nous définissons
le lecteur d'incidence minimal par
\[
 Q_{\min}=\operatorname{Supp}\circ c\circ\Psi:
 \mathcal U_{468}\longrightarrow\mathcal P([12]).
\]

Pour une application \(f:X\to Y\), notons
\[
 \operatorname{Eq}(f)=\{(x,x')\in X^2:f(x)=f(x')\}
\]
sa relation d'équivalence par fibres.

\begin{theorem}[Obstruction du lecteur exceptionnel minimal]
\label{p12-83-c17-s4:thm:obstruccion-lector-excepcional-minimo}
Dans le catalogue complet de 468 états, on vérifie
\[
 \operatorname{Eq}(\Psi,Q_{\min})=\operatorname{Eq}(\Psi).
\]
Dans la notation abrégée habituelle des partitions par fibres,
\[
 \ker(\Psi,Q_{\min})=\ker\Psi.
\]
Par conséquent, \((\Psi,Q_{\min})\) n'est pas injective et le lecteur
d'incidence minimal ne distingue aucun couple d'états déjà
identifié par \(\Psi\).
\end{theorem}

\begin{proof}
La factorisation \(Q_{\min}=S\circ\Psi\), avec
\(S=\operatorname{Supp}\circ c\), implique
\[
 \Psi(u)=\Psi(v)\Longrightarrow Q_{\min}(u)=Q_{\min}(v).
\]
L'implication réciproque entre les relations d'équivalence est immédiate,
puisque \(\Psi\) est la première composante de l'application diagonale.
L'égalité des relations s'obtient donc sans hypothèses statistiques.
L'énumération exhaustive confirme en outre que les deux partitions possèdent le
même histogramme :
\[
 132\,[1]+66\,[2]+18\,[3]+3\,[4]+6\,[5]+18\,[6],
\]
où \(a\,[m]\) signifie qu'il existe \(a\) fibres de cardinal \(m\).
\end{proof}

La microfibre de \(w_\pi=010211\) fournit le témoin distingué. Ses
cinq régions sont
\[
\begin{split}
 &(501,614,498,169,272,272),\quad
 (810,923,870,169,272,272),\\
 &(810,983,810,169,272,272),\quad
 (870,923,810,169,272,272),\\
 &(870,923,810,418,674,521).
\end{split}
\]
Les cinq ont la même image visible et le même support exceptionnel
\[
 P_\pi=\{2,4,5,6,7,8,9,10,11\}.
\]
Au niveau des 243 mots, on trouve 213 supports distincts : 185 ont
un antécédent, 27 en ont deux et un en a quatre. La prise de support introduit
donc un second quotient en plus de celui imposé par \(\Psi\).

Le théorème délimite le lecteur minimal, non l'évaluation exceptionnelle
enrichie. Un lecteur capable de séparer les régions d'un même mot
doit agir avant le quotient visible et utiliser au moins l'une des données
que celui-ci élimine : orientation, feuillet, signature, retour ou état de frontière.
L'entrelaceur dodécaphasique développé dans la partie suivante appartient à cette
classe enrichie et ne contredit pas l'obstruction précédente.

L'énumération complète des fibres produit l'histogramme précédent et
démontre que la prise de support introduit un second quotient.


\section{Publication archimédienne et d'incidence}

Pour chaque niveau, soit \(\mathcal E_n\) un objet d'incidence et soient
\begin{equation}
 e_n:S_n\longrightarrow\mathcal E_n,
 \qquad
 \sigma_{n,m}:\mathcal E_n\longrightarrow\mathcal E_m
 \label{p12-83-c17-s5:eq:datos-lector-incidencial-u024}
\end{equation}
des applications qui satisfont
\begin{equation}
 \sigma_{n,m}\circ e_n=e_m\circ r_{n,m}.
 \label{p12-83-c17-s5:eq:naturalidad-lector-incidencial-u024}
\end{equation}
La propriété universelle construit
\begin{equation}
 \operatorname{ev}_{\rm inc}:\Omega_\infty
 \longrightarrow\mathcal E_\infty,
 \qquad
 \mathcal E_\infty:=\varprojlim_n\mathcal E_n.
 \label{p12-83-c17-s5:eq:evaluacion-incidencial-limite-u024}
\end{equation}
La double publication est
\begin{equation}
 \mathcal P_\infty
 :=(\operatorname{ev}_{\mathbb R},\operatorname{ev}_{\rm inc}):
 \Omega_\infty\longrightarrow\mathbb R\times\mathcal E_\infty.
 \label{p12-83-c17-s5:eq:doble-publicacion-continuo-u024}
\end{equation}

\begin{figure}[H]
\centering
\begin{tikzcd}[column sep=10em,row sep=5em,
  cells={nodes={font=\large}}]
&\Omega_\infty
 \arrow[dl,"\operatorname{ev}_{\mathbb R}"']
 \arrow[dr,"\operatorname{ev}_{\rm inc}"]&\\
\mathbb R
 \arrow[dr,"\operatorname{coord}"']&&
\mathcal E_\infty
 \arrow[dl,"\operatorname{pub}_{\rm inc}"]\\
&\mathcal O&
\end{tikzcd}
\caption{Deux publications d'une même histoire. La commutativité dans
\(\mathcal O\) n'identifie pas les codomaines \(\mathbb R\) et
\(\mathcal E_\infty\).}
\label{p12-83-c17-s5:fig:doble-publicacion-continuo-u024}
\end{figure}

La naturalité démontrée dans les
\cref{p12-83-c17-s2:thm:lector-excepcional-global-rev6,p12-83-c17-s2:thm:doble-realizacion-global-rev6}
s'exprime, après application des lecteurs de publication, au moyen de
\begin{equation}
 \operatorname{pub}_{\rm inc}\circ\operatorname{ev}_{\rm inc}
 =\operatorname{coord}\circ\operatorname{ev}_{\mathbb R},
 \label{p12-83-c17-s5:eq:compatibilidad-doble-publicacion-u024}
\end{equation}
et produit une comparaison typée des deux publications. L'identité
conserve l'origine commune sans construire une branche à partir de l'autre : toutes deux
partent de \(\Omega_\infty\).


\section{Jauge de la publication exceptionnelle}

La seconde projection n'est pas un lecteur unaire dépourvu de données
d'incidence. Soit
\[
 \mathscr C_{\rm exc}
 =\mathscr C_{\rm ord}\times\mathscr C_{\rm Paley}
  \times\mathscr C_{\rm dod}\times\mathscr C_{\rm inc}
  \times\mathscr C_{\rm ret}
\]
l'espace des jauges exceptionnelles, dont les coordonnées fixent
respectivement l'ordre projectif, la jauge de Paley, la dodécade,
l'alignement d'incidence et la chambre de réseau. On fixe
\[
 \chi_{\rm exc}
 =(\chi_{\rm ord},\chi_{\rm Paley},\chi_{\rm dod},
   \chi_{\rm inc},\chi_{\rm ret})\in\mathscr C_{\rm exc}.
\]
Le lecteur typé a pour domaine
\[
 \mathfrak E:
 X_\infty^{\rm cal}\times\mathscr C_{\rm exc}
 \longrightarrow\mathbf{Exc},
 \qquad
 \mathfrak E_{\chi_{\rm exc}}(u):=\mathfrak E(u,\chi_{\rm exc}).
\]
Désormais, toute abréviation \(\mathfrak E(u)\) signifie
\(\mathfrak E_{\chi_{\rm exc}}(u)\) pour cette jauge explicitement fixée.
La dépendance à l'égard de \(\chi_{\rm exc}\) n'affecte pas l'existence du
continu ni sa publication archimédienne : elle type la conservation exceptionnelle
de l'incidence une fois l'état commun construit.

\begin{falsifier}[Perte de la jauge exceptionnelle]
La comparaison demeure incomplète si l'on change l'ordre projectif, la dodécade,
l'alignement d'incidence ou la chambre de réseau sans transporter les autres
coordonnées de la jauge, ou si l'on utilise \(\mathfrak E\) comme lecteur unaire avant
de fixer \(\chi_{\rm exc}\). Une telle défaillance invalide cette comparaison exceptionnelle ;
elle ne modifie pas le quotient archimédien et ne régénère pas le continu.
\end{falsifier}


La seconde lecture ne naît pas du nombre réel. Soit
\[
 \mathfrak E:
 \mathscr O_{\mathbb R,\leftrightarrow}^{\rm HMT}
 \longrightarrow\mathbf{Exc}
\]
le lecteur continu d'incidence, à codomaine séparé, qui prolonge
la chaîne de Paley--Witt, Golay,
Leech, \(V^\natural\), Monstre et Moonshine. Les deux applications partent de
la même histoire :
\begin{equation}
 \mathfrak D(u)=(\mathfrak R(u),\mathfrak E(u)).
 \label{p12-83-c17-s7:hmtch:eq-doble-proyeccion-two-way}
\end{equation}
La première contracte les fibres compatibles jusqu'à leur moyenne archimédienne ; la
seconde conserve incidence, orientation et mémoire de bord.

L'origine du calendrier dodécaphasique étant fixée, soit
\[
 R_+:=\{k\in\mathbb Z:1+(k\bmod 9)\in\{1,4,7\}\}
\]
l'ensemble bi-infini des retours complets \((\tau=+1)\). Il est syndétique :
ses lacunes sont uniformément bornées. Pour
\(u=(m,(x_k)_{k\in\mathbb Z})\), nous définissons
\[
 \operatorname{Ret}_+(u)
 :=\bigl(m,(x_k)_{k\in R_+}\bigr)
 \in\mathbb Z\times\prod_{k\in R_+}X,
\]
mais chaque \(x_k\) est ici l'état entier
\eqref{p12-83-c19-s1:hmtch:eq-celula-trit-completa} avec feuillet, retenue, route et incidence,
non le signe isolé.

\begin{proposition}[Échantillonnage sans perte du registre généalogique de retour]
\label{p12-83-c17-s7:hmtch:prop-retorno-reconstruye}
L'application \(\operatorname{Ret}_+\), munie de la topologie induite dans le
produit, est un homéomorphisme sur son image.
En conséquence, il existe une unique application continue
\(\mathfrak E_+\) sur cette image telle que
\begin{equation}
 \mathfrak E=\mathfrak E_+\circ\operatorname{Ret}_+.
 \label{p12-83-c17-s7:hmtch:eq-factorizacion-excepcional-retorno}
\end{equation}
\end{proposition}

\begin{proof}
Entre deux retours consécutifs, il y a un nombre uniformément borné de pas
du calendrier. L'état complet lors d'un retour contient la transition et
le registre généalogique qui déterminent les pas intermédiaires ; dans l'autre direction, on
utilise la flèche inverse de la complétion en groupoïde. Toute
l'orbite est ainsi reconstruite. L'application est continue et, lorsqu'on la restreint à une carte
compacte, une injection continue dans un espace séparé est un
homéomorphisme sur son image. La coordonnée entière est conservée
explicitement. La factorisation est définie en transportant \(\mathfrak E\) par
l'inverse et est continue carte par carte.
\end{proof}

C'est pourquoi le \(+1\) participe au corridor exceptionnel comme mémoire complète
de retour : l'échantillonnage d'états entiers ne perd pas la sortie exceptionnelle.
Cette proposition démontre la récupération et la factorisation, non une génération
du corridor à partir du scalaire \(+1\) isolé. La sortie archimédienne et la
sortie exceptionnelle sont deux résultantes coordonnées du même antécédent.
